\documentclass[10pt]{article}
\usepackage[margin=1in]{geometry}
\usepackage[T1]{fontenc}
\usepackage{lmodern,microtype,amsmath,amssymb,amsthm}
\theoremstyle{definition}
\newtheorem{assumption}{Assumption}
\begin{document}
\section*{Euclidean Strong-Growth Stochastic Minimization}
We consider minimizing a differentiable convex function $f:\mathbb{R}^d\to\mathbb{R}$, where $d$ is a positive integer. A given point $x^0\in\mathbb{R}^d$ lies within distance $R>0$ of a global minimizer $x^\star$. The stochastic gradient is unbiased, is Lipschitz in mean square with constant $\bar{L}>0$, and has variance bounded by $\Gamma^2\|\nabla f(x)\|^2$ at each $x$, where $\Gamma\geq0$. Each oracle call returns gradients at two queried points using a shared sample, with independent, identically distributed samples across calls. For $\varepsilon>0$, the goal is a measurable output $\bar{x}\in\mathbb{R}^d$ with integrable objective gap satisfying $\mathbb{E}_{P}[f(\bar{x})-f(x^\star)]\leq\varepsilon$, where $P$ is the law of the algorithm's and oracle's randomness.

The optimization problem is
\[
\min_{x\in\mathbb{R}^d} f(x),
\]
where $\|\cdot\|$ denotes the Euclidean norm. The parameters satisfy
\[
\begin{aligned}
\bar{L}&>0, & \Gamma&\geq0, & R&>0, & \varepsilon&>0,\\
67108864\sqrt{1+\Gamma^2}\,\varepsilon
&\leq \bar{L}R^2,\\
\left(\frac{\bar{L}R^2}{\sqrt{1+\Gamma^2}\,\varepsilon}\right)^3
&\leq d.
\end{aligned}
\]
Let $\Omega$ be a measurable sample space containing distinct points $a,b$ whose singletons are measurable. The algorithm receives $d$, $\Omega$ with its measurable structure, $\bar{L}$, $\Gamma$, $R$, $\varepsilon$, and $x^0$; the objective and the oracle distribution are accessed through oracle replies.

Our assumptions:
\begin{assumption}[Convexity and initial distance]
The function $f$ is Fr\'echet differentiable everywhere. For every $x,y\in\mathbb{R}^d$ and $t\in[0,1]$,
\[
f((1-t)x+ty)\leq (1-t)f(x)+t f(y).
\]
There exists $x^\star\in\mathbb{R}^d$ such that
\[
\begin{aligned}
f(x^\star)&\leq f(x) &&\text{for every }x\in\mathbb{R}^d,\\
\|x^0-x^\star\|&\leq R.
\end{aligned}
\]
\end{assumption}

\begin{assumption}[Unbiasedness and strong growth]
There exist a probability measure $\nu$ on $\Omega$ and a jointly measurable map $g:\mathbb{R}^d\times\Omega\to\mathbb{R}^d$. For every $x\in\mathbb{R}^d$, the map $g(x,\cdot)$ is Bochner integrable, the squared error $\|g(x,\cdot)-\nabla f(x)\|^2$ is integrable, and
\[
\begin{aligned}
\mathbb{E}_{\xi\sim\nu}[g(x,\xi)]
&=\nabla f(x),\\
\mathbb{E}_{\xi\sim\nu}
 \bigl[\|g(x,\xi)-\nabla f(x)\|^2\bigr]
&\leq \Gamma^2\|\nabla f(x)\|^2.
\end{aligned}
\]
\end{assumption}

\begin{assumption}[Mean-square smoothness and paired sampling]
For every $x,y\in\mathbb{R}^d$, the squared difference $\|g(x,\cdot)-g(y,\cdot)\|^2$ is integrable and satisfies
\[
\mathbb{E}_{\xi\sim\nu}
 \bigl[\|g(x,\xi)-g(y,\xi)\|^2\bigr]
\leq \bar{L}^2\|x-y\|^2.
\]
The oracle uses a sequence $(\xi_k)_{k\geq0}$ of measurable, independent random variables, each with law $\nu$, independent of the algorithm's randomness. At call $k$, a query $(x_k,y_k)\in\mathbb{R}^d\times\mathbb{R}^d$ returns
\[
\bigl(g(x_k,\xi_k),\ g(y_k,\xi_k)\bigr).
\]
Both components use the same sample, and the pair counts as one oracle call.
\end{assumption}

Queries may be chosen adaptively anywhere in $\mathbb{R}^d\times\mathbb{R}^d$. After the allotted paired oracle calls, the output $\bar{x}$ must be measurable, its objective gap must be integrable under $P$, and it must satisfy
\[
\mathbb{E}_{P}\bigl[f(\bar{x})-f(x^\star)\bigr]
\leq\varepsilon.
\]
\end{document}
